\documentclass[11pt,a4paper]{article}
\usepackage[utf8]{inputenc}
\usepackage[margin=1in]{geometry}
\usepackage{amsmath}
\usepackage{amssymb}
\usepackage{enumitem}
\usepackage{booktabs}
\usepackage{titlesec}
\usepackage{microtype}
\usepackage{xcolor}

\titleformat{\section}{\large\bfseries\sffamily}{}{0em}{}[\titlerule]
\titleformat{\subsection}{\normalsize\bfseries\sffamily}{}{0em}{}

\setlength{\parindent}{0pt}
\setlength{\parskip}{8pt plus 2pt minus 1pt}

\title{\textbf{CAMS Exam: Missing Topics \& High-Frequency Areas}}
\author{\textbf{Based on ACAMS Candidate Feedback}}
\date{\today}

\begin{document}
	
	\maketitle
	
	\section*{IMPORTANT NOTE FROM CANDIDATES}
	Based on recent exam-taker feedback:
	\begin{itemize}
		\item \textbf{3/4 of questions require selecting 2-3 answers}
		\item \textbf{FATF and USA PATRIOT Act are the most heavily tested topics}
		\item \textbf{Law enforcement cooperation (MLATs, Egmont, subpoenas) appears frequently}
		\item \textbf{Don't assume real-world experience replaces studying the ACAMS framework}
	\end{itemize}
	
	\section*{Instructions}
	Each question has \textbf{five options (A through E)}. 
	\begin{itemize}
		\item Questions marked \textbf{[SELECT TWO]} have exactly two correct answers
		\item Questions marked \textbf{[SELECT THREE]} have exactly three correct answers
	\end{itemize}
	
	\newpage
	
	% ============================================================
	\section{FATF Mutual Evaluations and Immediate Outcomes}
	% ============================================================
	
	\subsection{1. FATF Mutual Evaluation: Two Pillars \textnormal{[SELECT TWO]}}
	
	A FATF member nation is undergoing its Mutual Evaluation. The assessment team notes that the country has enacted all 40 Recommendations into national law, has a fully operational FIU receiving 50,000+ STRs annually, and has implemented CDD and wire transfer rules. However, the country has: (i) zero money laundering convictions in five years; (ii) never frozen terrorist assets; (iii) never conducted a financial investigation on its own initiative; and (iv) no record of international cooperation requests. The country's legal framework is perfect on paper.
	
	Which two statements correctly describe how this country will be rated?
	
	\begin{enumerate}[label=\Alph*.]
		\item The country will receive a rating of Compliant or Largely Compliant on \textbf{Technical Compliance} because the laws exist on paper mapping to all 40 Recommendations.
		\item The country will receive a rating of Low on \textbf{Effectiveness} across multiple Immediate Outcomes, particularly IO.6 (financial intelligence used), IO.7 (ML investigations and prosecutions), IO.8 (confiscation), and IO.10 (TF investigations and prosecutions).
		\item A country cannot receive a low effectiveness rating if it has perfect technical compliance with all Recommendations.
		\item The Mutual Evaluation process only assesses technical compliance, not effectiveness in practice.
		\item FATF Mutual Evaluations are conducted annually for all member countries.
	\end{enumerate}
	
	\textbf{Correct Answers: A and B}
	
	\subsection{2. FATF Immediate Outcomes (IO.1-IO.11) \textnormal{[SELECT THREE]}}
	
	A FATF assessor is preparing to evaluate a country's effectiveness. The assessor must understand which outcomes measure specific aspects of the AML/CFT system.
	
	Which three of the following correctly match an Immediate Outcome to its description?
	
	\begin{enumerate}[label=\Alph*.]
		\item \textbf{IO.1} — ML/TF risks are understood and, where appropriate, actions coordinated domestically.
		\item \textbf{IO.6} — Financial intelligence is used by competent authorities for ML/TF investigations.
		\item \textbf{IO.8} — Proceeds of crime and instrumentalities of crime are confiscated.
		\item \textbf{IO.2} — International cooperation is used to obtain evidence from foreign jurisdictions (this is actually IO.2, not IO.1).
		\item \textbf{IO.11} — Financial institutions and DNFBPs are supervised for compliance with AML/CFT requirements.
	\end{enumerate}
	
	\textbf{Correct Answers: A, B, and C}
	
	\subsection{3. FATF Grey List vs. Black List: Institutional Response \textnormal{[SELECT THREE]}}
	
	A multinational bank has correspondent relationships with banks in two jurisdictions: Country Alpha (recently placed on FATF Grey List — Increased Monitoring) and Country Beta (on FATF Black List — Call for Action). The bank's compliance committee must determine the appropriate response to each.
	
	Which three of the following correctly describe the bank's obligations?
	
	\begin{enumerate}[label=\Alph*.]
		\item For Country Alpha (Grey List), the bank must apply Enhanced Due Diligence to all transactions and correspondent accounts but is not automatically required to terminate relationships.
		\item For Country Beta (Black List), FATF's call for countermeasures means the bank should evaluate and implement appropriate countermeasures, which may include severely restricting or terminating transactions and correspondent relationships.
		\item Both Grey List and Black List jurisdictions require identical treatment: all transactions must be frozen and all accounts closed immediately.
		\item A Grey List jurisdiction has committed to resolving identified AML deficiencies within agreed timeframes, which may justify continued relationships with enhanced monitoring.
		\item Countermeasures against Black List jurisdictions are optional for private financial institutions; only central banks are obligated to act.
	\end{enumerate}
	
	\textbf{Correct Answers: A, B, and D}
	
	\subsection{4. FATF Recommendation 16: Wire Transfer Rule \textnormal{[SELECT TWO]}}
	
	A payments processor is auditing its compliance with FATF Recommendation 16 (the Travel Rule). The audit reviews 500 cross-border wire transfers.
	
	Which two statements correctly describe the required information for transfers above USD/EUR 1,000?
	
	\begin{enumerate}[label=\Alph*.]
		\item Originator information must include: name, account number, and at least one of — address, national ID number, customer identification number, or date and place of birth.
		\item Beneficiary information must include: name and account number.
		\item For transfers below USD/EUR 1,000, no originator or beneficiary information is required.
		\item Only the originator's account number is required; name is optional for domestic transfers.
		\item The Travel Rule applies only to fiat currency transfers; virtual asset transfers are exempt.
	\end{enumerate}
	
	\textbf{Correct Answers: A and B}
	
	\subsection{5. FATF Recommendation 3: Predicate Offenses \textnormal{[SELECT TWO]}}
	
	A country is drafting its money laundering legislation and must determine which predicate offenses to include.
	
	Which two statements correctly describe FATF's requirements for predicate offenses?
	
	\begin{enumerate}[label=\Alph*.]
		\item FATF recommends the "all-crimes approach," under which money laundering applies to proceeds from all criminal acts defined as serious under national law.
		\item FATF mandates a minimum of 22 designated predicate offenses, including drug trafficking, fraud, corruption, and tax crimes.
		\item Money laundering should apply to persons who committed the predicate offense (self-laundering is criminalized).
		\item Predicate offenses are limited to drug trafficking and terrorism financing only.
		\item Self-laundering is explicitly excluded from criminalization under FATF Recommendation 3.
	\end{enumerate}
	
	\textbf{Correct Answers: A and C}
	
	% ============================================================
	\section{USA PATRIOT Act (Heavily Tested)}
	% ============================================================
	
	\subsection{6. USA PATRIOT Act Section 311: Designation and Measures \textnormal{[SELECT TWO]}}
	
	The U.S. Secretary of the Treasury has designated a foreign financial institution as a "primary money laundering concern" under Section 311 of the USA PATRIOT Act.
	
	Which two measures can be imposed on U.S. financial institutions regarding this foreign institution?
	
	\begin{enumerate}[label=\Alph*.]
		\item Prohibition or strict conditioning of the opening or maintaining of correspondent or payable-through accounts for that foreign institution.
		\item Imposition of special recordkeeping and reporting requirements on transactions involving the designated institution.
		\item Automatic seizure of all assets held by the foreign institution in U.S. banks.
		\item Criminal indictment of all officers of the foreign institution.
		\item Mandatory closure of all U.S. bank branches in the foreign institution's home country.
	\end{enumerate}
	
	\textbf{Correct Answers: A and B}
	
	\subsection{7. USA PATRIOT Act Section 313: Shell Bank Prohibition \textnormal{[SELECT TWO]}}
	
	A U.S. correspondent bank discovers that one of its foreign respondent banks has no physical presence in any country, employs no staff, and is not affiliated with a regulated depository institution. The respondent bank processes USD 200 million annually through the correspondent account.
	
	Which two statements correctly describe the correspondent bank's obligations?
	
	\begin{enumerate}[label=\Alph*.]
		\item A "shell bank" is defined as a bank with no physical presence in any country and not affiliated with a regulated financial institution.
		\item U.S. banks are absolutely prohibited from maintaining correspondent accounts for foreign shell banks under Section 313.
		\item Shell banks are permitted as long as they are regulated by a G20 central bank.
		\item The correspondent bank may continue the relationship if it files a SAR every 90 days.
		\item An exception exists for shell banks that are affiliates of U.S.-regulated banks.
	\end{enumerate}
	
	\textbf{Correct Answers: A and B}
	
	\subsection{8. USA PATRIOT Act Section 314(a): Law Enforcement Requests \textnormal{[SELECT TWO]}}
	
	A financial institution receives a FinCEN 314(a) request identifying 15 individuals suspected of laundering proceeds from a transnational drug trafficking organization.
	
	Which two statements correctly describe the institution's obligations?
	
	\begin{enumerate}[label=\Alph*.]
		\item The institution must search its accounts and transactions according to FinCEN's instructions for accounts maintained within the time period specified.
		\item The institution must maintain the confidentiality of the request and cannot disclose it to the subjects of the request.
		\item The institution must automatically freeze any matching account without further investigation.
		\item The institution may ignore the request if the individuals are not current customers.
		\item The institution must file a SAR for every account searched regardless of whether a match is found.
	\end{enumerate}
	
	\textbf{Correct Answers: A and B}
	
	\subsection{9. USA PATRIOT Act Section 314(b): Information Sharing \textnormal{[SELECT TWO]}}
	
	Two banks, NorthBank and SouthBank, each have identified suspicious activity involving accounts held by different individuals who appear to be part of the same money laundering network. Neither bank has enough information alone to build a complete case.
	
	Which two statements correctly describe Section 314(b) information sharing?
	
	\begin{enumerate}[label=\Alph*.]
		\item Under Section 314(b), financial institutions that have filed the required FinCEN opt-in notification may voluntarily share information with other opted-in institutions about persons they suspect of ML/TF activity.
		\item The sharing institutions receive a safe harbor from liability for the sharing.
		\item Section 314(b) requires both institutions to first obtain a court order before exchanging any customer information.
		\item Section 314(b) sharing automatically constitutes tipping-off because the shared information relates to a suspected subject of an AML investigation.
		\item Section 314(b) sharing is only permissible if the amount of suspected money laundering exceeds \$1 million.
	\end{enumerate}
	
	\textbf{Correct Answers: A and B}
	
	\subsection{10. USA PATRIOT Act Section 319(b): Record Production and Account Seizure \textnormal{[SELECT TWO]}}
	
	A U.S. bank maintains a correspondent account for a foreign bank. Federal law enforcement serves a subpoena on the U.S. bank demanding records related to the foreign bank's account. The foreign bank refuses to comply.
	
	Which two statements correctly describe Section 319(b) of the USA PATRIOT Act?
	
	\begin{enumerate}[label=\Alph*.]
		\item Section 319(b) allows the U.S. government to demand that a foreign bank produce records related to its U.S. correspondent account.
		\item If the foreign bank refuses to comply with a subpoena, the U.S. bank may be required to terminate the correspondent account, and funds in the account may be seized.
		\item Section 319(b) only applies to banks located in FATF Black List jurisdictions.
		\item Foreign banks are never required to produce records to U.S. authorities.
		\item Section 319(b) requires the U.S. bank to pay the foreign bank's legal fees.
	\end{enumerate}
	
	\textbf{Correct Answers: A and B}
	
	% ============================================================
	\section{Law Enforcement Cooperation (MLATs, Egmont, Subpoenas)}
	% ============================================================
	
	\subsection{11. Mutual Legal Assistance Treaties (MLATs) \textnormal{[SELECT TWO]}}
	
	Country A has identified that \$20 million in embezzled government funds has been deposited into a bank account in Country B. Country A and Country B have an active Mutual Legal Assistance Treaty (MLAT). Country A's prosecutor needs to obtain the bank records to support a criminal prosecution.
	
	Which two statements correctly describe the function of an MLAT?
	
	\begin{enumerate}[label=\Alph*.]
		\item MLATs provide a formal, legally binding mechanism for governments to cooperate in the exchange of evidence, judicial documents, and asset tracking/freezing across borders.
		\item An MLAT request is the proper channel for obtaining bank records that will be admissible in criminal proceedings in Country A's courts.
		\item MLATs allow Country A to send its police forces to arrest the bank manager in Country B without local authorization.
		\item Country A's FIU can obtain the same records through an Egmont Group request without an MLAT, and those records will be admissible in criminal court.
		\item MLATs are informal, non-binding agreements with no legal force.
	\end{enumerate}
	
	\textbf{Correct Answers: A and B}
	
	\subsection{12. Egmont Group: FIU-to-FIU Information Exchange \textnormal{[SELECT TWO]}}
	
	An investigator at Financial Intelligence Unit (FIU) Alpha has identified a suspicious cross-border wire sequence that terminates in the jurisdiction of FIU Beta. The investigator needs additional information to determine whether a crime has occurred.
	
	Which two statements correctly describe FIU-to-FIU information exchange under the Egmont Group framework?
	
	\begin{enumerate}[label=\Alph*.]
		\item FIUs can exchange financial intelligence directly using the Egmont Group's secure communication platform (Egmont Secure Web).
		\item Exchanged information must be used only for AML/CFT purposes and the receiving FIU must maintain confidentiality.
		\item Egmont Group information exchange can replace an MLAT for criminal prosecution purposes (intelligence from Egmont is automatically admissible in court).
		\item Any FIU in the world can join the Egmont Group without any membership requirements.
		\item Information exchanged through the Egmont Group can be sold to commercial data brokers.
	\end{enumerate}
	
	\textbf{Correct Answers: A and B}
	
	\subsection{13. Responding to Law Enforcement Subpoenas \textnormal{[SELECT TWO]}}
	
	A bank receives a formal federal law enforcement subpoena demanding the immediate production of all account records, wire transfer logs, and signature cards for a specific corporate customer under investigation for money laundering.
	
	Which two actions are appropriate for the compliance officer?
	
	\begin{enumerate}[label=\Alph*.]
		\item Promptly comply with the subpoena within the specified legal timeframe.
		\item Maintain a complete internal mirror file of everything produced to law enforcement.
		\item Notify the customer immediately so they can clear the account before production.
		\item Refuse to provide the documentation citing bank-client confidentiality.
		\item Destroy all responsive records to protect customer privacy.
	\end{enumerate}
	
	\textbf{Correct Answers: A and B}
	
	\subsection{14. SAR Tipping-Off: Permissible vs. Prohibited Disclosures \textnormal{[SELECT TWO]}}
	
	A bank has filed an STR on a customer. Three scenarios occur: (i) the relationship manager tells the customer "compliance is reviewing some transactions on your account"; (ii) the compliance team shares the STR narrative with the bank's legal department for internal legal review; (iii) a second bank contacts the compliance department as part of a joint 314(b) investigation and the compliance team shares its STR findings.
	
	Which two statements correctly identify tipping-off violations?
	
	\begin{enumerate}[label=\Alph*.]
		\item Scenario (i) constitutes tipping-off — telling a customer they are under compliance review, even without mentioning the STR by name, violates AML confidentiality laws.
		\item Scenario (ii) does NOT constitute tipping-off — sharing an STR with the bank's own legal department on a need-to-know basis is permissible.
		\item Scenario (iii) constitutes tipping-off — sharing STR information with another bank is never permissible.
		\item All three scenarios constitute tipping-off because any disclosure of STR information for any purpose is prohibited.
		\item Scenario (i) is permissible because the relationship manager did not mention the STR by name.
	\end{enumerate}
	
	\textbf{Correct Answers: A and B}
	
	\subsection{15. Law Enforcement "Keep Open" Requests Post-SAR \textnormal{[SELECT TWO]}}
	
	A bank has filed three SARs on an account over six months. The compliance committee decides to terminate the relationship and close the account. Before closing, a federal law enforcement agent contacts the bank in writing, requesting that the bank keep the account open for an additional 90 days to allow an ongoing undercover investigation to continue.
	
	Which two protocols should the bank follow?
	
	\begin{enumerate}[label=\Alph*.]
		\item Document the request in writing from law enforcement before agreeing to keep the account open.
		\item Obtain senior management approval before agreeing to the "keep open" request.
		\item Close the account immediately regardless of law enforcement's request to minimize regulatory risk.
		\item Publicly announce the surveillance operation to deter criminal activity on the account.
		\item Ignore the request because the bank has already decided to close the account.
	\end{enumerate}
	
	\textbf{Correct Answers: A and B}
	
	% ============================================================
	\section{BSA Recordkeeping, Safe Harbor, and Whistleblower Protections}
	% ============================================================
	
	\subsection{16. SAR Safe Harbor and Confidentiality \textnormal{[SELECT TWO]}}
	
	A bank files a SAR on a customer. The customer later sues the bank for defamation. An employee accidentally mentions the SAR filing to a third party during a social conversation.
	
	Which two statements correctly describe SAR Safe Harbor and confidentiality rules?
	
	\begin{enumerate}[label=\Alph*.]
		\item The SAR Safe Harbor provides immunity from civil liability to financial institutions that file SARs in good faith.
		\item Unauthorized disclosure of a SAR filing is a criminal offense under 31 U.S.C. \S 5318(g)(2), punishable by fines and imprisonment.
		\item A SAR and its supporting documentation must be retained for three years from the date of filing.
		\item The SAR Safe Harbor does NOT protect the institution if the SAR was filed in bad faith or with reckless disregard for the truth.
		\item An institution may use a SAR as evidence to defend itself in a civil lawsuit brought by the customer.
	\end{enumerate}
	
	\textbf{Correct Answers: A and B}
	
	\subsection{17. BSA Recordkeeping Requirements (5-Year Rule) \textnormal{[SELECT THREE]}}
	
	A compliance director is reviewing document retention policies. The bank maintains the following five categories of records.
	
	Which three categories must be retained for FIVE YEARS under the BSA?
	
	\begin{enumerate}[label=\Alph*.]
		\item CTR filings and all supporting documentation.
		\item SAR filings and all supporting documentation.
		\item Wire transfer records for transactions of \$3,000 or more (31 CFR \S 1020.410).
		\item Customer identification program (CIP) documents — may be retained for only two years after account closure.
		\item Daily branch cash logs — no retention requirement under the BSA.
	\end{enumerate}
	
	\textbf{Correct Answers: A, B, and C}
	
	\subsection{18. Whistleblower Protections Under the BSA \textnormal{[SELECT TWO]}}
	
	A compliance analyst reports internally that senior management is suppressing SAR filings on a high-revenue corporate client. Three weeks later, the analyst receives a performance improvement plan (PIP) — her first negative review in eight years. She believes the PIP is retaliatory.
	
	Which two statements correctly describe whistleblower protections available to her?
	
	\begin{enumerate}[label=\Alph*.]
		\item The Bank Secrecy Act (31 U.S.C. \S 5328) specifically prohibits financial institutions from retaliating against employees who report BSA violations.
		\item The analyst may file a complaint with the U.S. Department of Labor (OSHA), which investigates BSA whistleblower retaliation claims, within 180 days of the adverse employment action.
		\item The analyst is protected only if she filed the report externally with FinCEN; internal hotline reports are not covered.
		\item There is no federal whistleblower protection for AML compliance analysts; the bank's internal HR policy is the only recourse.
		\item The analyst must first obtain a lawyer and file a federal lawsuit within 30 days of the retaliation.
	\end{enumerate}
	
	\textbf{Correct Answers: A and B}
	
	% ============================================================
	\section{BSA/AML Program Governance and Independent Testing}
	% ============================================================
	
	\subsection{19. BSA/AML Program: Five Required Pillars \textnormal{[SELECT THREE]}}
	
	Under the Bank Secrecy Act (31 CFR \S 1020.210) as amended by FinCEN's 2016 CDD Final Rule, a financial institution's AML compliance program must include five minimum pillars.
	
	Which three of the following are among those five pillars?
	
	\begin{enumerate}[label=\Alph*.]
		\item Internal controls reasonably designed to assure compliance with the BSA.
		\item Independent testing for compliance conducted by bank personnel or an outside party.
		\item A designated compliance officer responsible for managing BSA compliance.
		\item Quarterly certification to the board of directors that all SARs have been filed.
		\item Annual profit-sharing with regulatory examiners.
	\end{enumerate}
	
	\textbf{Correct Answers: A, B, and C}
	
	\subsection{20. Three Lines of Defense Model \textnormal{[SELECT TWO]}}
	
	A bank's internal audit department is asked to help redesign the AML transaction monitoring scenarios because the compliance team lacks technical expertise. The audit team agrees and co-designs 14 new scenarios. Six months later, the same audit team conducts an independent test of the system.
	
	Which two statements correctly identify the governance failure?
	
	\begin{enumerate}[label=\Alph*.]
		\item The internal audit team has compromised its independence by designing the same controls it subsequently audited. This violates the Three Lines of Defense framework.
		\item The compliance team (Second Line) operates the AML program; internal audit (Third Line) must remain independent and cannot design operational controls.
		\item The compliance team acted appropriately by delegating technical design work to internal audit; this is standard industry practice.
		\item The governance violation can be cured by having the Chief Compliance Officer sign an attestation confirming the audit findings are independent.
		\item Internal audit reporting to the Audit Committee is sufficient to maintain independence regardless of its involvement in design.
	\end{enumerate}
	
	\textbf{Correct Answers: A and B}
	
	\subsection{21. Independent Testing: Qualification and Reporting \textnormal{[SELECT TWO]}}
	
	A bank's audit committee is selecting an individual or team to perform the independent testing of the AML program.
	
	Which two characteristics are required for the independent testing function?
	
	\begin{enumerate}[label=\Alph*.]
		\item The tester must be completely independent of the functions being tested.
		\item The tester must possess sufficient expertise and qualifications in AML/CFT.
		\item The tester must report directly to the Chief Compliance Officer.
		\item The tester must be a current employee of the bank (external auditors are not permitted).
		\item The tester must be an active law enforcement officer.
	\end{enumerate}
	
	\textbf{Correct Answers: A and B}
	
	% ============================================================
	\section{Risk-Based Approach and Enterprise Risk Assessment}
	% ============================================================
	
	\subsection{22. Risk-Based Approach (RBA) vs. Rules-Based Approach \textnormal{[SELECT TWO]}}
	
	A compliance officer is explaining the risk-based approach to a new board member who asks, "Why don't we just apply the same due diligence measures to every customer? That would be simpler."
	
	Which two statements correctly justify the risk-based approach under FATF standards?
	
	\begin{enumerate}[label=\Alph*.]
		\item An RBA concentrates resources and controls on areas of highest risk, making the overall system more effective and efficient.
		\item Proportionate measures are applied to lower-risk situations without creating unnecessary burdens on customers or the institution.
		\item The RBA eliminates all due diligence requirements for low-risk customers.
		\item Regulators prefer rules-based approaches because they are easier to audit.
		\item The RBA is only applicable to correspondent banking relationships, not retail customers.
	\end{enumerate}
	
	\textbf{Correct Answers: A and B}
	
	\subsection{23. Enterprise-Wide Risk Assessment (EWRA) \textnormal{[SELECT TWO]}}
	
	A bank is preparing its annual Enterprise-Wide Risk Assessment. The risk team has identified inherent risk scores for each business line and mitigating controls.
	
	Which two components must be documented in the EWRA?
	
	\begin{enumerate}[label=\Alph*.]
		\item An assessment of inherent risk exposure across all products, services, customer types, geographies, and delivery channels.
		\item The mitigating controls in place and the resulting net residual risk determination.
		\item A list of all SARs filed in the previous calendar year sorted by dollar amount.
		\item The personal tax returns of the board of directors.
		\item The bank's quarterly profit margin by business line.
	\end{enumerate}
	
	\textbf{Correct Answers: A and B}
	
	% ============================================================
	\section{Beneficial Ownership and PEPs}
	% ============================================================
	
	\subsection{24. Beneficial Ownership: Aggregating Direct and Indirect Stakes \textnormal{[SELECT ONE]}}
	
	A compliance officer is onboarding a new client, Apex Industrial Corp. The ownership structure is:
	\begin{itemize}
		\item Individual K owns 20\% of Apex directly.
		\item Individual K also owns 80\% of Zenith Holdings LLC, which owns 15\% of Apex.
	\end{itemize}
	The regulatory UBO threshold is 25\% combined direct and indirect ownership.
	
	What is Individual K's total effective ownership stake in Apex, and what is the correct compliance treatment?
	
	\begin{enumerate}[label=\Alph*.]
		\item 20\% (direct only) — below threshold — no verification required.
		\item 32\% (20\% direct + (80\% × 15\%)) — exceeds threshold — must be identified, verified, and documented as a UBO.
		\item 35\% (20\% + 15\% pass-through) — exceeds threshold — UBO verification required.
		\item Individual K does not qualify as a UBO because the indirect stake is held through a separate legal entity.
		\item The 25\% threshold applies only to direct ownership; indirect stakes are irrelevant.
	\end{enumerate}
	
	\textbf{Correct Answer: B}
	
	\subsection{25. Politically Exposed Persons (PEPs) Classification \textnormal{[SELECT TWO]}}
	
	A private bank is onboarding a new client, Ms. Y, who is the adult daughter of the Minister of Energy of an oil-rich nation. Ms. Y does not hold any government position herself. She wishes to deposit \$15 million from "consulting fees paid by energy logistics companies."
	
	Which two statements correctly describe her PEP classification and required due diligence?
	
	\begin{enumerate}[label=\Alph*.]
		\item Ms. Y should be classified as a PEP family member, requiring Enhanced Due Diligence (EDD) and senior management approval.
		\item EDD for PEP family members requires establishing source of wealth and source of funds through verifiable, corroborative documentation.
		\item Ms. Y is not a PEP because she holds no government position herself, so standard CDD is sufficient.
		\item The ambiguous "consulting fees" from energy firms related to her father's ministry is not a red flag.
		\item Only the Minister of Energy himself is a PEP; family members are never classified as PEPs.
	\end{enumerate}
	
	\textbf{Correct Answers: A and B}
	
	% ============================================================
	\section{Correspondent Banking and Wolfsberg}
	% ============================================================
	
	\subsection{26. Wolfsberg Correspondent Banking Due Diligence Questionnaire (CBDDQ) \textnormal{[SELECT THREE]}}
	
	A correspondent bank is reviewing a CBDDQ submitted by a prospective respondent bank.
	
	Which three areas are typically covered in the Wolfsberg CBDDQ?
	
	\begin{enumerate}[label=\Alph*.]
		\item The respondent bank's AML/CFT policies, procedures, and internal controls.
		\item The respondent bank's customer base, including high-risk segments like MSBs, NBFIs, and foreign PEPs.
		\item The respondent bank's independent testing schedule and results.
		\item The respondent bank's CEO's personal investment portfolio.
		\item The respondent bank's office furniture preferences.
	\end{enumerate}
	
	\textbf{Correct Answers: A, B, and C}
	
	\subsection{27. Nested Correspondent Accounts and Undisclosed Downstream Clients \textnormal{[SELECT TWO]}}
	
	PrimeClear Bank (EU) provides a USD correspondent account to RegionalBank, a commercial bank in a jurisdiction with weak AML controls. PrimeClear discovers that RegionalBank has been routing transactions for CryptoX, an unregulated virtual asset exchange, through the correspondent account without disclosure. CryptoX was never approved by PrimeClear.
	
	Which two statements correctly describe PrimeClear's risks and obligations?
	
	\begin{enumerate}[label=\Alph*.]
		\item PrimeClear has lost visibility into the true originators and beneficiaries of the nested transactions, which violates core correspondent banking due diligence standards.
		\item PrimeClear must issue a Request for Information to RegionalBank, file an STR if warranted, and reassess whether the relationship should continue.
		\item PrimeClear has no liability because RegionalBank holds a valid domestic banking license.
		\item Nested transactions are a routine feature of correspondent banking and require no additional due diligence.
		\item PrimeClear should immediately terminate the relationship without any prior investigation.
	\end{enumerate}
	
	\textbf{Correct Answers: A and B}
	
	% ============================================================
	\section{De-Risking and Financial Inclusion}
	% ============================================================
	
	\subsection{28. Indiscriminate De-Risking Consequences \textnormal{[SELECT TWO]}}
	
	A global bank terminates all correspondent relationships with banks headquartered in a region based solely on the region's FATF rating, without conducting individual risk assessments of the 18 respondent banks. Several of those banks have been rated Compliant by their local regulators and have passed recent FATF mutual evaluations.
	
	Which two statements correctly describe the consequences of this blanket termination?
	
	\begin{enumerate}[label=\Alph*.]
		\item Legitimate financial institutions with robust AML programs lose access to USD clearing, causing financial exclusion and driving transactions into unregulated, informal channels.
		\item The blanket termination violates the risk-based approach, which requires individual, documented risk assessments rather than categorical exclusions based solely on geography.
		\item The bank has acted properly because banks have absolute discretion to choose their correspondent partners.
		\item The termination will be praised by regulators as a model of de-risking.
		\item The bank faces no consequences because de-risking is explicitly required by FATF.
	\end{enumerate}
	
	\textbf{Correct Answers: A and B}
	
	% ============================================================
	\section{CDD, EDD, and SDD Eligibility}
	% ============================================================
	
	\subsection{29. Simplified Due Diligence (SDD): Appropriate Application \textnormal{[SELECT TWO]}}
	
	A compliance officer is reviewing five account relationships to determine whether Simplified Due Diligence (SDD) is appropriate.
	
	Which two types of customers may qualify for SDD under a risk-based approach?
	
	\begin{enumerate}[label=\Alph*.]
		\item A publicly traded company listed on a recognized stock exchange with full SEC disclosure.
		\item A U.S. government agency (e.g., Social Security Administration).
		\item A foreign PEP from a high-corruption jurisdiction.
		\item A shell company registered in a secrecy haven.
		\item A customer who refuses to provide any identification documents.
	\end{enumerate}
	
	\textbf{Correct Answers: A and B}
	
	\subsection{30. Enhanced Due Diligence (EDD): Source of Wealth vs. Source of Funds \textnormal{[SELECT ONE]}}
	
	A private banker is conducting EDD on a PEP client, the Minister of Natural Resources of a West African country. The minister wishes to deposit \$4.5 million from the sale of a Paris property. He provides a notarized sale contract. When asked about source of wealth (how he accumulated his overall wealth), he states he has been a civil servant and inherited property, but declines to provide documentation.
	
	Why is the banker's proposed approach (accepting only the sale contract) insufficient?
	
	\begin{enumerate}[label=\Alph*.]
		\item The sale contract is not an acceptable source of funds document; only bank statements qualify.
		\item EDD for PEPs requires establishing BOTH source of funds (the specific transaction) AND source of wealth (overall asset accumulation). The unresolved question — how a lifelong civil servant acquired a \$4.5 million Paris property — must be answered through corroborative documentation.
		\item The minister's refusal to provide documentation is a legal right and cannot be used as a red flag.
		\item EDD for a government minister is only required if the minister's country is on the FATF Grey List.
	\end{enumerate}
	
	\textbf{Correct Answer: B}
	
	% ============================================================
	\section{Conclusion — Final High-Frequency Topics}
	% ============================================================
	
	\subsection{31. AML Training: Role-Based Requirements \textnormal{[SELECT TWO]}}
	
	A bank's external auditor rated the AML training program as "ineffective" because the same training module was delivered to all employees, from the board of directors to front-line tellers.
	
	Which two principles should guide the redesigned training program?
	
	\begin{enumerate}[label=\Alph*.]
		\item Training must be tailored to the specific roles, responsibilities, and ML/TF risks faced by each employee group.
		\item Board members need training on governance oversight and ML/TF risk strategy, while tellers need training on cash structuring indicators and CTR reporting.
		\item All employees must receive identical training to demonstrate consistent institutional coverage.
		\item Training is only required for employees who have direct customer contact; back-office staff have no training obligations.
		\item The IT infrastructure team has no material exposure to AML risk and can be exempted.
	\end{enumerate}
	
	\textbf{Correct Answers: A and B}
	
	\subsection{32. Adverse Media and OSINT: Escalation Triggers \textnormal{[SELECT TWO]}}
	
	A bank's EDD team discovers through OSINT that a corporate customer's CEO was named (not charged) in a foreign parliamentary corruption inquiry, a predecessor company with the same address was dissolved after a trade fraud investigation, and the customer's primary counterparty's director appears on an Interpol Red Notice.
	
	Which two statements correctly describe the required response?
	
	\begin{enumerate}[label=\Alph*.]
		\item The OSINT findings collectively constitute credible adverse media requiring immediate escalation to the compliance committee.
		\item OSINT findings do not require a conviction to be actionable — reasonable suspicion of ML/TF risk, not proof of guilt, is the standard.
		\item Because no conviction has occurred and the account has been clean for four years, the adverse media is inconclusive.
		\item The Interpol Red Notice alone is sufficient to close the account immediately without further investigation.
		\item The dissolved predecessor company finding is irrelevant because it was a separate legal entity.
	\end{enumerate}
	
	\textbf{Correct Answers: A and B}
	
\end{document}